\section{Spectral locality and the temperature multiplier field}
  \label{sec:spectral-locality}

  \subsection{Local spectral energy density}
    \label{subsec:local-density}

    Let $\Ag$ be a positive Laplace-type operator on $(M,g)$.
    Denote by $\hk(x,y;t)$ its heat kernel, satisfying
    \begin{equation}
      (\partial_t + \Ag)\hk(x,y;t) = 0,
      \qquad
      \hk(x,y;0^+) = \delta_g(x,y).
      \label{eq:heat-eq}
    \end{equation}
    The diagonal $u_0(x;t) \equiv \hk(x,x;t)$ is a local scalar density encoding
    the spectral weight at point $x$ and diffusion scale $t$.

    \begin{definition}[Local spectral energy density]
      \label{def:u}
      The \emph{local spectral energy density} at point $x$ and scale $t$ is
      \begin{equation}
        u(x;t) \equiv -\partial_t \log \hk(x,x;t).
        \label{eq:u-def}
      \end{equation}
    \end{definition}

    This definition is minimal: it depends only on local diagonal heat-kernel data.
    The sign convention ensures $u(x;t) > 0$ in the short-time regime, consistent with
    an energy interpretation.
    The expression is also the natural local analogue of the canonical relation
    $E = -\partial_\beta \log Z$ in statistical mechanics, once $t$ is identified with
    an inverse spectral resolution parameter.

  \subsection{Seeley--DeWitt expansion of \texorpdfstring{$u(x;t)$}{u(x;t)}}
    \label{subsec:SDW-expansion}

    For the minimal scalar Laplacian $\Ag = -\nabla_g^2$ in dimension $d=4$, the
    short-time expansion of the diagonal heat kernel reads~\cite{Vassilevich2003,
  Gilkey1995}
    \begin{equation}
      \hk(x,x;t) \sim \frac{1}{(4\pi t)^2}
      \bigl(a_0(x) + a_2(x)\,t + a_4(x)\,t^2 + \cdots\bigr)
      \quad\text{as }t\to 0^+,
      \label{eq:hk-expansion}
    \end{equation}
    with Seeley--DeWitt coefficients
    \begin{align}
      a_0(x) &= 1, \label{eq:a0}\\
      a_2(x) &= \tfrac{1}{6}R(x), \label{eq:a2}\\
      a_4(x) &= \frac{1}{(4\pi)^2\cdot 360}
      \bigl(R_{\mu\nu\rho\sigma}R^{\mu\nu\rho\sigma}
      - R_{\mu\nu}R^{\mu\nu} + \tfrac{5}{3}R^2\bigr). \label{eq:a4}
    \end{align}

    Taking the logarithm of~\eqref{eq:hk-expansion} and differentiating with respect
    to $t$ yields:
    \begin{equation}
      \partial_t \log \hk(x,x;t)
      = -\frac{2}{t}
      + \partial_t \log\bigl(a_0 + a_2\,t + a_4\,t^2 + \cdots\bigr).
      \label{eq:dlog-hk}
    \end{equation}
    For the minimal Laplacian, $a_0 = 1$ and $a_2 = \frac{1}{6}R$, so
    \begin{equation}
      \partial_t \log\bigl(1 + \tfrac{1}{6}R\,t + O(t^2)\bigr)
      = \tfrac{1}{6}R + O(t).
      \label{eq:log-expand}
    \end{equation}
    Therefore, by Definition~\ref{def:u}:

    \begin{proposition}[Curvature expansion of $u$]
      \label{prop:u-expansion}
      For the minimal scalar Laplacian in dimension four,
      \begin{equation}
        u(x;t) = \frac{2}{t} - \frac{1}{6}R(x) + O(t).
        \label{eq:u-expansion}
      \end{equation}
      The leading term is universal and UV-divergent; the first geometric correction is
      proportional to the scalar curvature with coefficient fixed by $a_2$.
    \end{proposition}

  \subsection{The local multiplier field \texorpdfstring{$\beta(x)$}{beta(x)}}
    \label{subsec:beta-field}

    We now evaluate $u(x;t)$ at a fixed spectral scale $t = t_*$, defined as the
    admissibility scale introduced in~\cite{Gravity20,Spectral12} as the infrared
    threshold below which the spectral description is consistent.

    \begin{definition}[Local spectral multiplier field]
      \label{def:beta}
      The \emph{local spectral multiplier field} $\beta(x)^{-1}$ is defined by
      \begin{equation}
        \beta(x)^{-1} \equiv u(x;t_*) = \frac{2}{t_*} - \frac{1}{6}R(x) + O(t_*).
        \label{eq:beta-def}
      \end{equation}
    \end{definition}

    \begin{remark}[Sign convention and decomposition]
      \label{rem:sign}
      Equation~\eqref{eq:beta-def} may appear to give a \emph{negative} curvature correction,
      whereas the literature on local temperature often writes a \emph{positive} geometric
      correction.
      The sign is consistent and arises from the definition $u = -\partial_t \log K$
      together with the expansion $\log K \sim -2\log t + \frac{1}{6}Rt + \cdots$,
      which gives $\partial_t \log K \sim -2/t + \frac{1}{6}R$,
      and hence $u \sim 2/t - \frac{1}{6}R$.

      When we write the decomposition with $\beta_0^{-1} \equiv 2/t_*$, we have
      \begin{equation}
        \beta(x)^{-1} = \beta_0^{-1} - \frac{1}{6}R(x) + O(t_*).
        \label{eq:beta-decomp}
      \end{equation}
      On a positively curved manifold ($R > 0$), the local multiplier is \emph{smaller}
      than the flat reference $\beta_0^{-1}$: curvature reduces the effective spectral
      temperature relative to flat space.
      This sign is physically consistent with the Seeley--DeWitt interpretation:
      positive curvature lowers the local spectral density, hence lowers the effective
      temperature.
    \end{remark}

    \begin{remark}[Status of $t_*$ and $\beta_0^{-1}$]
      \label{rem:tstar}
      The scale $t_*$ is a spectral admissibility parameter, not a physical temperature.
      Its value is fixed by the spectral resolution of the relational system, as discussed
      in~\cite{Spectral12}.
      The universal part $\beta_0^{-1} = 2/t_*$ is a UV quantity reflecting the spectral
      density of states at the cutoff.
      It is scheme-dependent and carries no intrinsic thermodynamic meaning at this stage.
      Only the geometric correction $\frac{1}{6}R(x)$ is covariant and physically
      significant: it encodes local curvature as a deviation from the flat-space multiplier.
    \end{remark}

    \begin{remark}[Decomposition structure]
      \label{rem:decomp}
      The decomposition~\eqref{eq:beta-decomp} has the form
      \[
        \beta(x)^{-1}
        = \underbrace{\beta_0^{-1}}_{\text{universal (UV)}}
        - \underbrace{\tfrac{1}{6}R(x)}_{\text{geometric correction}}.
      \]
      The geometric correction carries a negative sign: positive curvature reduces the
      local spectral multiplier relative to the flat-space reference $\beta_0^{-1}$.
      The correction is entirely controlled by the $a_2$ Seeley--DeWitt coefficient,
      which is precisely the coefficient that generates the Einstein--Hilbert counterterm
      via pole subtraction in the companion paper.
      The two papers are therefore connected at the level of the same spectral coefficient.
    \end{remark}
